\usepackage[utf8]{inputenc}

\ifarxiv
\usepackage[a4paper,left=2cm,right=2cm,top=2cm,bottom=2cm]{geometry}
\usepackage[onehalfspacing]{setspace}
\usepackage[numbers,longnamesfirst]{natbib}
\usepackage{amsthm,amsthm,amssymb}
\usepackage{authblk}
\else
\citestyle{acmauthoryear} % required by ACM journals
\fi

\usepackage{amsmath,mathtools} % amssymb?
\usepackage{dsfont}
%\usepackage{upgreek}

\usepackage{booktabs} % spacing and styling
\usepackage{xspace}
\usepackage{graphicx}
\usepackage{xspace}
\usepackage{url}
%\usepackage{enumerate}
\usepackage{enumitem}
\setlist[enumerate,1]{label=\textup{(\roman*)}}

\usepackage{caption} % captions ans subfigures
\usepackage{subcaption}

\usepackage[algo2e,ruled,vlined,linesnumbered]{algorithm2e} % algorithm type setting
\SetAlgoSkip{}

%\usepackage{balance} % balancing the last page

\allowdisplaybreaks[3] % break equations between pages

%% -- drawing and plotting library support

\usepackage{tikz}
\usetikzlibrary{shapes,arrows}
\usetikzlibrary{decorations}
\usetikzlibrary{plotmarks}
\usetikzlibrary{calc}
\usetikzlibrary{mindmap}
\usetikzlibrary{shadows}
\usetikzlibrary{backgrounds}
\usetikzlibrary{patterns}
\usetikzlibrary{shapes.symbols}

\usepackage{pgfplots}
\usepackage{pgfplotstable}
\usepgfplotslibrary{statistics}
%\usetikzlibrary{external}
%\tikzexternalize[prefix=figures/]
\pgfplotsset{compat=1.18}

%% -- customise the theorem environments

%\newtheorem{theorem}{Theorem}
\newtheorem{algorithm}{Algorithm}
%\newtheorem{remark}[theorem]{Remark}
\newtheorem{remark}{Remark}
%\newtheorem{proposition}[theorem]{Proposition}
\newtheorem{openproblem}{Open Problem}

\ifarxiv
\newtheorem{theorem}             {Theorem}
\newtheorem{lemma}      [theorem]{Lemma}
\newtheorem{corollary}  [theorem]{Corollary}
\newtheorem{conjecture} [theorem]{Conjecture}
\newtheorem{definition} [theorem]{Definition}
\newtheorem{proposition}[theorem]{Proposition}
\fi

\usepackage{tabularx,colortbl}
\newcommand{\myc}{\cellcolor{orange!30}}

\makeatletter
\g@addto@macro\bfseries{\boldmath}
\makeatother


%% -- English/Latin abbreviations

\newcommand{\ie}{i.\,e.\xspace}
\newcommand{\wrt}{w.\,r.\,t.\xspace}
\newcommand{\eg}{e.\,g.\xspace}
\newcommand{\whp}{w.\,h.\,p.\xspace}
\newcommand{\Wlog}{W.\,l.\,o.\,g.\xspace}
\newcommand{\wlo}{w.\,l.\,o.\,g.\xspace}

%% -- mathematical symbols and notations

\newcommand{\Z}{\mathbb{Z}}
\newcommand{\N}{\mathbb{N}}
\newcommand{\R}{\mathbb{R}}

\newcommand{\filt}{\mathcal{F}}
\newcommand{\filtt}{\mathcal{F}_t}

\DeclareMathOperator{\XOR}{\oplus} % XOR
\DeclareMathOperator*{\hdist}{H}   % Hamming distance

\DeclareMathOperator*{\argmax}{argmax}
\DeclareMathOperator*{\argmin}{argmin}

\DeclareMathOperator{\Unif}{Unif}                         % uniform distribution
\DeclareMathOperator{\Bin}{Bin}                           % binomial distribution
\DeclareMathOperator{\Poi}{Poi}                           % Poisson distribution
\DeclareMathOperator{\Nor}{\ensuremath{\mathcal{N}}}      % Normal distribution

\newcommand*{\smallO}{\mathrm{o}}
\newcommand*{\bigO}{\mathcal{O}}
\newcommand*{\eulerE}{\mathrm{e}}

\DeclareMathOperator{\poly}{poly}                         % polynomial

\DeclareMathOperator*{\convhull}{ConvHull}

\newcommand{\prob}[1]{\Pr\left(#1\right)}                 % probability
\newcommand{\Prob}[1]{\Pr\left(#1\right)}                 % probability
\newcommand{\expect}[1]{\mathrm{E}\left[#1\right]}        % expectation
\newcommand{\E}[1]{\mathrm{E}\left[#1\right]}        % expectation
\newcommand{\var}[1]{\mathbf{Var}\left[#1\right]}         % variance

\newcommand{\ind}[1]{\mathds{1}_{#1}}                     % indicator of an event
\newcommand{\indc}[1]{\mathds{1}_{\{#1\}}}                % indicator of a condition

\newcommand{\xmin}{x_{\min}}
\newcommand{\xmax}{x_{\max}}

\newcommand{\vecone}{\vec{1}}

\newcommand{\ndom}[2]{\ensuremath{\mathrm{ND}_{#1}\left(#2\right)}} % non-dominated set
\newcommand{\Popt}[1]{\ensuremath{\mathcal{P}_{#1}}}                % set of all Pareto optimal points
\newcommand{\ineff}[2]{(#1,#2)-Pareto-sparse}                       % (eps,l)-sparsity
\newcommand{\ineffg}[2]{(#1,#2)-Pareto-isolated}                         % (eps,l)-gap


%\newcommand{\filtcond}[1]{\,;\,{#1}\mid\filt_t}
%\newcommand{\filtzero}{\mathcal{F}_0}

%\newcommand*{\potentialFunction}{g}
%\newcommand*{\potentialSum}{Y}
%\newcommand*{\differenceInPotential}{\Delta}

%\newcommand{\card}[1]{\left|#1\right|}
%\newcommand{\indic}[1]{\mathds{1}\{#1\}}

%%% -- function, operator, and algorithm names
\newcommand{\RRRMO}{\ensuremath{\textsc{RR}_{\mathrm{MO}}}\xspace}  % the Real Royal Road function
\newcommand{\lRRRMO}{\ensuremath{\textsc{RRR-MO}'}\xspace}% the local version without the global front

\newcommand{\posFit}{\ensuremath{G}\xspace}               % set of strictly positive fitness (f1>0,f2>0)
\newcommand{\frontP}{\ensuremath{F}\xspace}               % Pareto optimal set of RRRMO
\newcommand{\lfrontP}{\ensuremath{F'}\xspace}             % Pareto optimal set of lRRRMO

\newcommand{\hv}{\ensuremath{\mathrm{HV}}\xspace}         % hypervolume
\newcommand{\deltahv}{\ensuremath{\mathrm{HV}^+}\xspace}  % contribution of hypervolume

\newcommand{\nadir}{\ensuremath{\mathrm{nad}}\xspace}      % nadir point
\newcommand{\ideal}{\ensuremath{\min}\xspace}              % ideal point

\newcommand{\tree}{\mathcal{T}}

\newcommand{\LZ}{\textsc{LZ}\xspace}                      % # of leading zeroes
\newcommand{\TZ}{\textsc{TZ}\xspace}                      % # of trailing zeroes
\newcommand{\LO}{\textsc{LO}\xspace}                      % # of leading ones

\newcommand{\ones}[1]{|#1|_1}                             % # of ones
\newcommand{\zeroes}[1]{|#1|_0}                           % # of zeroes
\newcommand{\OM}{\textsc{OneMax}\xspace}                           % OneMax function
\newcommand{\ZM}{\textsc{ZeroMax}\xspace}                          % ZeroMax function

\newcommand{\JUMP}{\textsc{Jump}\xspace}                           % Jump function
\newcommand{\TRAP}{\textsc{Trap}\xspace}                           % Trap function
\newcommand{\ONEMAX}{\textsc{OneMax}\xspace}                       % OneMax function
\newcommand{\TWOMAX}{\textsc{TwoMax}\xspace}                       % TwoMax function
\newcommand{\SAT}{\textsc{Sat}\xspace}                             % Sat decision problem
\newcommand{\MAXSAT}{\textsc{MaxSat}\xspace}                       % MaxSat problem
\newcommand{\COUNTSAT}{\textsc{CountSat}\xspace}                   % CountSat function
\newcommand{\OZBinval}{\textsc{OZSoftBinVal}\xspace}
\newcommand{\Binval}{\textsc{BinVal}\xspace}
\newcommand{\BV}{\textsc{OZBinVal}\xspace}


\newcommand{\RRRfull}{\textsc{RealRoyalRoad}\xspace}               % RRR full name

\newcommand{\OZCSAT}{\textsc{OZCSat}\xspace}               % OZCOUNTSAT function
\newcommand{\BICOUNTSAT}{\textsc{OZCSat}\xspace}               % another name 
\newcommand{\OZCSATfull}{\textsc{OneZeroCountSat}\xspace}      % OZCOUNTSAT full name

\newcommand{\COCZ}{\textsc{COCZ}\xspace}                           % COCZ function
\newcommand{\COCZfull}{\textsc{CountingOnesCountingZeroes}\xspace} % COCZ full name

\newcommand{\OMM}{\textsc{OMM}\xspace}                             % OMM function
\newcommand{\OMMfull}{\textsc{OneMinMax}\xspace}                   % OMM full name

\newcommand{\LOTZ}{\textsc{LOTZ}\xspace}                           % LOTZ function
\newcommand{\LOTZfull}{\textsc{LeadingOnesTrailingZeroes}\xspace}  % LOTZ full name

\newcommand{\OJZJ}{\textsc{OJZJ}\xspace}                           % OJZJ function
\newcommand{\OJZJfull}{\textsc{OneJumpZeroJump}\xspace}            % OJZJ full name

\newcommand{\OTZT}{\textsc{OTZT}\xspace}                           % OTZT function
\newcommand{\OTZTfull}{\textsc{OneTrapZeroTrap}\xspace}            % OTZT full name

\newcommand{\ZOmonotone}{0/1-monotone\xspace}                      % 0/1-monotone
\newcommand{\LOLZmonotone}{LO/LZ-monotone\xspace}                  % LO/TZ-monotone

\newcommand{\cdist}{\ensuremath{\textsc{cDist}}\xspace}   % crowding distance
\DeclareMathOperator{\BinTour}{\textsc{BinaryTournament}} % binary tournament

\newcommand{\muplusga}{($\mu$+$\lambda$)~GA\xspace}       % (mu+lambda) GA algorithm

\newcommand{\Pdouble}{P^{\mathrm{double}}}
\newcommand{\pot}[1]{\varphi_{#1}}
\newcommand{\multisetminus}{\mathrel{-}}
%\newcommand{\COUNTSAT}{\textsc{CountSat}\xspace}
%\newcommand{\BICOUNTSAT}{\textsc{OneZeroCountSat}\xspace}
%\newcommand{\MAXSAT}{\textsc{MaxSat}\xspace}
%\newcommand{\SAT}{\textsc{Sat}\xspace}

\newcommand{\nsga}{NSGA\nobreakdash-II\xspace}
\newcommand{\nsgaiii}{NSGA\nobreakdash-III\xspace}
\newcommand{\refer}{\mathcal{R}}

\newcommand{\alggsemo}{GSEMO\xspace}
\newcommand{\algnsga}{NSGA\nobreakdash-II\xspace}
\newcommand{\algnsgaiiidef}{NSGA\nobreakdash-III$_{\mathrm{def}}$\xspace}
\newcommand{\algsmsdef}{SMS-EMOA$_{\mathrm{def}}$\xspace}
\newcommand{\algnsgaiii}{NSGA\nobreakdash-III\xspace}
\newcommand{\algsms}{SMS-EMOA\xspace}

% squish some space
\ifarxiv\else 
\setlength{\textfloatsep}{1ex}
\setlength{\floatsep}{1ex}
\setlength{\intextsep}{2ex}
%\setlength{\dbltextfloatsep}{0pt}
%\setlength{\dblfloatsep}{0pt}
\setlength{\abovecaptionskip}{1ex}
\setlength{\belowcaptionskip}{0ex}
\fi

%% -- other macros

%use the following command to switch the plots on (first line) or off (second line)
%\newcommand{\toggleplot}[1]{#1}
\newcommand{\toggleplot}[1]{{\textcolor{red}{Plots removed to increase compilation speed. Use command $\backslash$toggleplot in preamble to reinsert them.}}}

\iffinal
\newcommand{\andre}[1]{\textcolor{green!50!black}{}}
\newcommand{\dirk}[1]{\textcolor{purple}{}}
\newcommand{\cuong}[1]{\textcolor{orange}{}}
\newcommand{\newedit}[1]{\textcolor{black}{#1}}
\newcommand{\neweditx}[1]{\textcolor{black}{#1}}
\newcommand*{\todo}[1]{\textcolor{red}{}}
\else
\newcommand{\andre}[1]{\textcolor{green!50!black}{[(Andre) #1]}}
\newcommand{\dirk}[1]{\textcolor{purple}{[(Dirk) #1]}}
\newcommand{\cuong}[1]{\textcolor{orange}{[(Cuong) #1]}}
\newcommand{\newedit}[1]{\textcolor{blue}{#1}}
\newcommand{\neweditx}[1]{\textcolor{blue!20!white}{#1}}
\newcommand*{\todo}[1]{\textcolor{red}{\textrm{(TODO: #1)}}}
\fi
